\section{El dilema de la evaluación en el entrenamiento de modelos}

\subsection{Fallo de los Benchmarks y generalización del Agent}

\begin{warningbox}{Riesgo de sobreajuste en los Benchmarks}
El campo de los Agent enfrenta actualmente un problema serio de evaluación:
\begin{itemize}
    \item La cantidad de Benchmarks de Agent disponibles es limitada
    \item Elevar la puntuación de un Benchmark no representa la capacidad real: muchas veces se trata de un sobreajuste a una tarea específica
    \item La percepción de los usuarios en escenarios OOD (Out-of-Distribution) se desliga de la puntuación del Benchmark
    \item «Quien siembra melones cosecha melones, quien siembra frijoles cosecha frijoles»: se mejora aquello para lo que se entrena, pero la generalización es limitada
\end{itemize}
\end{warningbox}

\subsection{El origen del dilema de evaluación}

Yang Zhilin (杨植麟) considera que el origen del dilema de evaluación es el siguiente:

\begin{enumerate}
    \item El aprendizaje por refuerzo depende en última instancia de una buena evaluación o verificación: los problemas matemáticos y la programación con casos de prueba (Test Case) pueden verificarse de manera automática, pero para las tareas complejas de extremo a extremo es difícil encontrar una buena forma de evaluación
    \item Los nuevos problemas que trae la mejora de la capacidad de los modelos: la complejidad de las tareas aumenta y la forma de evaluación necesita mantenerse al día
    \item La característica de distribución de cola larga de las tareas de Agent: es difícil cubrir todos los escenarios reales
\end{enumerate}

\subsection{Resumen de la sección}

El fallo de los Benchmarks y la generalización insuficiente de los Agent son el dilema de evaluación más destacado en la actualidad. El origen está en que las tareas complejas son difíciles de verificar de manera automática, así como en el Gap entre la distribución de entrenamiento y los escenarios reales de uso. Hacer que la IA participe en la evaluación y el entrenamiento (la IA entrena a la IA) podría ser una dirección de avance.

